% !TeX program = pdflatex
% !BIB program = bibtex
\PassOptionsToPackage{english}{babel}
\documentclass[final,1p,times,authoryear]{elsarticle}
\makeatletter
\def\ps@pprintTitle{%
  \let\@oddhead\@empty
  \let\@evenhead\@empty
  \let\@oddfoot\@empty
  \let\@evenfoot\@empty
}
\makeatother
\usepackage[utf8]{inputenc}
\usepackage[T1]{fontenc}
\usepackage{amsmath}
\usepackage{amsthm}
\usepackage{amssymb}
\usepackage{newtxtext,newtxmath}
\usepackage[english]{babel}
\usepackage[babel=true]{csquotes}
\usepackage{tcolorbox}
\usepackage{booktabs}
\usepackage{multirow}
\usepackage{float}
\usepackage[unicode=true,pdfencoding=auto]{hyperref}

\usepackage{threeparttable}
\usepackage{tikz}
\usepackage{array, longtable, tabularx}
\usepackage{adjustbox}
\usepackage{ltxtable}
\usepackage{makecell}
\usepackage{pdflscape}
\usepackage{enumitem}
\usepackage{wrapfig}
\usepackage[expansion=false]{microtype}
\usepackage{mathtools}
\usepackage{bookmark}
\usepackage[english]{cleveref}

\allowdisplaybreaks
\tolerance=1000
\emergencystretch=10pt
\hyphenpenalty=50
\exhyphenpenalty=50

\setlength{\textfloatsep}{10pt plus 2pt minus 2pt}
\setlength{\floatsep}{8pt plus 2pt minus 2pt}
\setlength{\intextsep}{8pt plus 2pt minus 2pt}
\setlength{\abovecaptionskip}{6pt}
\setlength{\belowcaptionskip}{6pt}
\setlength{\columnsep}{20pt}
\setlength{\parskip}{6pt plus 1pt minus 1pt}
\setlength{\parindent}{0pt}

\providecommand{\X}{X}
\providecommand{\T}{T}
\providecommand{\DD}{\mathrm{DD}}
\providecommand{\DI}{\mathrm{DI}}
\providecommand{\DT}{\mathrm{DT}}
\providecommand{\Lmax}{L_{\max}}
\providecommand{\epsedge}{\epsilon_{\text{edge}}}
\providecommand{\epscontrib}{\epsilon_{\text{contrib}}}

\begin{document}

\begin{frontmatter}
\title{Methodology of Elcano's Index on Economic Dependence (EIED): Critical Networks, Indirect Dependencies, and Structural Power in Global Trade}

\author[upo,elcano]{Manuel Hidalgo-Pérez\corref{cor1}}\ead{mhidalgo@rielcano.org}
\cortext[cor1]{Corresponding author: Manuel Hidalgo-Pérez, mhidalgo@rielcano.org}

\affiliation[upo]{
  organization={Universidad Pablo de Olavide},
  addressline={Ctra Utrera s/n},
  postcode={41013},
  city={Seville},
  country={Spain}
}

\affiliation[elcano]{
  organization={Real Instituto Elcano Foundation},
  addressline={C. de Príncipe de Vergara, 51},
  postcode={28006},
  city={Madrid},
  country={Spain}
}

\begin{abstract}
The rising geopolitical fragmentation has transformed the international trade landscape, exposing previously underestimated economic vulnerabilities. This paper proposes a new economic security indicator that overcomes the limitations of traditional approaches by systematically capturing indirect trade dependencies transmitted through intermediary countries. By integrating Input-Output matrices, network theory, and shock propagation algorithms, our indicator quantifies both direct and indirect dependencies, revealing hidden vulnerabilities in global supply chains. The empirical analysis demonstrates that traditional approaches significantly underestimate the exposure to external disruptions in strategic sectors. Domestic supply, missing for a growing share of cells in the underlying trade database, is completed with OECD input-output data through validated imputation rules, so that dependence is always measured against total supply. These findings have important implications for strategic autonomy policies, suggesting that effective diversification must consider the entire structure of the global trade network rather than being limited to direct bilateral relations.
\end{abstract}

\begin{keyword}
Economic security \sep Trade dependencies \sep Network theory \sep Strategic autonomy \sep Global value chains\\
\textbf{JEL Codes:} F14, F15, F52, C67, D85
\end{keyword}
\end{frontmatter}

\section{Introduction}

In a context of growing geoecononomic fragmentation, systemic tensions between major powers, and frequent disruptions in global supply chains, the evaluation of national economic security requires more refined analytical tools. Existing literature has predominantly focused on aggregate bilateral trade metrics or direct dependency, failing to capture the structural interdependencies that emerge through global production networks and indirect trade.

This paper addresses this gap by developing a Elcano's Index on Economic Dependence (EIED), designed to measure a country's exposure to exogenous disruptions by integrating direct and indirect dependencies, and considering the structural relevance of critical intermediaries.

For the purposes of this research, we define economic security as "the capacity of a nation to maintain the integrity of its essential economic functions and its critical supply chains in the face of exogenous disruptions, whether they arise from natural disasters, technological failures, or deliberate actions of other actors."

This definition distinguishes economic security from related concepts:
\begin{itemize}
    \item \textbf{Economic sovereignty:} Unlike sovereignty, which often implies autarky or minimal reliance on foreign inputs, economic security recognizes the reality and benefits of interdependence while focusing on managing critical vulnerabilities.
    \item \textbf{Resilience:} While resilience refers to the capacity to recover from disruptions, economic security encompasses both the prevention of critical disruptions and the maintenance of functionality during them.
    \item \textbf{Competitiveness:} In contrast with competitiveness, which focuses on relative advantage in global markets, economic security focuses on the stability and reliability of economic functions essential for national well-being.
    \item \textbf{Traditional national security:} Unlike traditional conceptions of national security, which primarily address military threats, economic security recognizes that existential threats can arise from supply disruptions, technological dependency, or economic coercion.
\end{itemize}

Our Elcano's Index on Economic Dependence operationalizes this definition by measuring both direct and indirect vulnerabilities across all industries, with special attention to sectors where disruptions would have cascading effects throughout the economy. The index is designed to identify not only obvious dependencies in direct bilateral trade but also hidden vulnerabilities arising from a country's position within complex global supply networks.

\section{Theoretical Foundations of the Elcano's Index on Economic Dependence}

This work anchors its approach to measuring economic security in three complementary theoretical frameworks: the theory of complex interdependence, structural power in international political economy, and weaponized interdependence.

\textbf{Complex Interdependence and Network Vulnerability}
Drawing on the foundational work of \cite{keohane1977power} on complex interdependence, we conceptualize economic security not simply as the absence of bilateral dependencies, but as an emergent property of a country's position within multiple overlapping networks of exchange. As \cite{farrell2019weaponized} point out, "interdependence has created complex networks with particular topographies," where certain nodes become critical bottlenecks for global flows. Our EIED operationalizes this network-based understanding by explicitly modeling how disruptions propagate through multiple paths within the global trade system.

\textbf{Structural Power in International Political Economy}
The second theoretical pillar is Strange's \citeyear{strange1988states} concept of structural power, understood as "the power to shape and determine the global political economy structures [...] within which other states, their political institutions, their economic enterprises [...] have to operate." In Strange's framework, structural power operates through control over security, production, finance, and knowledge. Our index specifically operationalizes structural power in the production domain by quantifying how control over critical nodes in supply chains confers leverage that extends far beyond direct bilateral relations. By identifying indirect dependencies and critical intermediaries, we make visible what Strange described as the "indirect and unobservable exercise of power" in global economic relations.

\textbf{Weaponized Interdependence}
Finally, our work engages with Farrell and Newman's \citeyear{farrell2019weaponized} theory of "weaponized interdependence," which posits that "states with political authority over the nodes in international networked structures through which money, information, and goods travel are uniquely positioned to impose costs on others." This theory distinguishes between "panopticon effects" (the ability to monitor flows) and "choke-points" (the ability to restrict flows). Our EIED specifically measures vulnerability to the latter.

\section{Motivation and Geoeconomic Context}

Economic security has emerged as a global strategic priority in a context marked by successive systemic crises that have revealed critical vulnerabilities in international supply chains. Recent years have shown how excessive dependency on specific suppliers can undermine not only economic stability but also national security in a broader sense.

The COVID-19 pandemic represented the first systemic shock that clearly exposed these fragilities. Disruptions in the production and distribution of essential goods demonstrated that geographic concentration of production capacities can lead to critical shortages with cascading effects across multiple economic sectors.

The Russian invasion of Ukraine further intensified these concerns, revealing particularly acute vulnerabilities in the European energy and food sectors. Recently, trade barriers have been used as diplomatic tools, highlighting how national security concerns are increasingly intertwined with trade policy. A recent example is the announcement of reciprocal tariffs by the US administration on imports from the European Union and China, underlining the vulnerability of trade relations to non-economic considerations.

In this context, access to critical raw materials, emerging technologies, and digital infrastructures has become central to strategic planning. Entities such as the European Union have developed "open strategic autonomy" strategies to balance the benefits of free trade with the need to secure access to essential resources.

\section{Literature Review}

\subsection{Evolution of Trade Dependency Measurement}
The measurement of trade dependency has been a central concern in international economics since the mid-20th century. Early approaches focused primarily on bilateral relations and import concentration. However, as global production fragmented across borders, the insufficiency of these traditional measures became apparent. The integration of global value chains (GVCs) has fundamentally altered the nature of trade dependencies.

\subsection{Methodological Approaches for Measuring Dependencies}
\subsubsection{Concentration Indices and Bilateral Measures}
One of the most widely used tools is the Herfindahl-Hirschman Index (HHI), which assesses the geographic concentration of imports. Institutions such as the European Commission have employed this index to identify strategic dependencies, establishing specific thresholds to signal risks to economic security. Nonetheless, these indices present limitations: they are static and fail to capture indirect dependencies arising through third countries.

\subsubsection{Input-Output Analysis and Intersectoral Linkages}
Input-Output (I-O) analysis, pioneered by \cite{Leontief1951}, provides a framework for understanding the network of intersectoral dependencies. The Leontief inverse matrix, $(I - A)^{-1}$, captures both the direct and indirect requirements of each sector to produce a unit of final demand. An important extension is the concept of "average propagation lengths" (APL) developed by \cite{Dietzenbacher2005}, which measures the economic distance between sectors through forward and backward linkages.

\subsubsection{Network Analysis and Centrality Measures}
Network analysis emerges as a complementary approach that models trade as a complex system of interconnected nodes (countries or sectors) and links (trade flows). Pioneers like \cite{Fagiolo2010} characterized the World Trade Network (WTN) documenting its structural properties. Centrality concepts such as betweenness allow the identification of nodes acting as structural "hubs" or critical bottlenecks.

\subsubsection{Shock Propagation Models}
Recent research focuses on the dynamic mechanisms through which economic shocks propagate. \cite{Acemoglu2012} and \cite{Carvalho2013} demonstrated that micro-level shocks to individual firms or sectors can have significant aggregate effects when production networks exhibit certain structural properties, such as a heavy-tailed degree distribution.

\section{Development of the Proposed Indicator}

\subsection{Conceptualization}
The proposal of this work is a comprehensive indicator that combines matrix analysis, network theory, and propagation algorithms to identify and quantify hidden vulnerabilities.

\subsubsection{Definition of Direct and Indirect Dependency}
Trade dependency is defined as the exposure of a country to supply disruptions originating in another, considering all possible channels.
\begin{enumerate}
    \item \textbf{Direct Dependency ($\DD$):} Proportion of direct imports in a country's total supply.
    \item \textbf{Indirect Dependency ($\DI$):} Proportion of supply attributable to an exporter through multi-step trade routes (intermediaries).
\end{enumerate}

\subsubsection{Advantages of the Proposed Indicator}
Our approach stands out for three key advantages: (1) identification of hidden systemic vulnerabilities that bilateral indices do not detect, (2) quantitative assessment of resilience to disruptions by modeling alternative routes, and (3) prospective analysis of geoecononomic fragmentation scenarios.

\begin{table}[!t]
\caption{Comparison of Methods for Measuring Trade Dependency} \label{tab:comparacion-metodos}
\centering
\footnotesize
\renewcommand{\arraystretch}{1.1}
\begin{adjustbox}{width=\textwidth,center}
    \begin{tabular}{@{}l*{6}{c}@{}}
        \toprule
        \makecell[l]{\textbf{Feature}} & \textbf{HHI} & \makecell{\textbf{Imp.}\\\textbf{Conc.}} & \textbf{TiVA} & \makecell{\textbf{Trad.}\\\textbf{I-O}} & \makecell{\textbf{Net-}\\\textbf{works}} & \makecell{\textbf{Prop.}\\\textbf{Ind.}} \\
        \midrule
        Captures direct dep. & \checkmark & \checkmark & \checkmark & \checkmark & \checkmark & \checkmark \\
        Captures indirect dep. & \texttimes & \texttimes & Partial & Direct only & \checkmark & \textbf{Complete} \\
        Normalisation (0-1) & \checkmark & \checkmark & \texttimes & \texttimes & \texttimes & \checkmark \\
        Input-Output based & \texttimes & \texttimes & \checkmark & \checkmark & \texttimes & \checkmark \\
        Network theory & \texttimes & \texttimes & \texttimes & \texttimes & \checkmark & \checkmark \\
        Models shock propagation & \texttimes & \texttimes & \texttimes & Limited & Partial & \checkmark \\
        Detects hidden vuln. & \texttimes & \texttimes & Partial & \texttimes & Partial & \checkmark \\
        Resilience analysis & \texttimes & \texttimes & Limited & Limited & Partial & \checkmark \\
        \bottomrule
    \end{tabular}
\end{adjustbox}
\end{table}

\subsection{Mathematical Formulation}
The total dependency $\text{DT}_{ij}$ of an importer $j$ on an exporter $i$ is defined as the sum of direct and indirect dependency.
We consider the normalized trade matrix $\mathbf{T}$, where each element $T_{ab} = x_{ab} / S_b$ represents the fraction of $b$'s supply that comes directly from $a$, and $S_b$ is $b$'s total supply of the industry: domestic production sold in $b$ plus imports (\cref{sec:data}).
The strength of a trade path $p = (i, k_1, \ldots, k_{\ell-1}, j)$ of length $\ell$ is calculated as the product of the dependencies along its edges:
\begin{equation}
F(p) = \prod_{(a,b) \in \text{edges}(p)} T_{ab} = T_{ik_1} \cdot T_{k_1k_2} \cdots T_{k_{\ell-1}j}
\end{equation}

The total dependency is the sum of the strengths of all economically significant paths from $i$ to $j$:
\begin{equation}
\text{DT}_{ij} = \sum_{\ell=1}^{L_{\max}} \sum_{p \in \mathcal{P}_{ij}^{\ell}} F(p)
\end{equation}

\subsection{Vulnerability Index and National Aggregation}
To convert the network of trade dependencies into an operational measure of economic security, the model calculates the \textbf{Vulnerability Index} at the industry and country levels. This process consists of two phases:

\begin{enumerate}
    \item \textbf{Industrial Vulnerability:} For each industry-country pair, the total dependency obtained from the network analysis is weighted by two risk factors:
    \begin{itemize}
        \item \textbf{Strategic criticality of the product:} Measures the difficulty of local substitution or the essentiality of the input for the functioning of the economy.
        \item \textbf{Geopolitical risk of the supplier:} Measured through the voting affinity index in the United Nations General Assembly (UNGA). Lower political affinity with the supplier increases the importer's vulnerability.
    \end{itemize}
    
    \item \textbf{Aggregation and National Vulnerability:} A country's National Vulnerability Index is obtained as the weighted sum of sector-level vulnerabilities, using the economic weight (the total value of imports in dollars) of each industry on the national total as weights:
    \begin{equation}
    V_j = \sum_{s} w_{js} \cdot V_{js}
    \end{equation}
    where $V_{js}$ is the vulnerability of country $j$ in sector $s$, and $w_{js}$ is the share of sector $s$ in $j$'s input imports. The final index is normalized in a range from 0 to 1, where 1 represents the maximum relative systemic vulnerability.
\end{enumerate}

\subsection{Computational Implementation and Practical Considerations}
The practical application to real data poses significant computational challenges due to the exponential growth of potential paths. Our implementation incorporates optimization strategies:
\begin{itemize}
    \item \textbf{Commercial relevance threshold ($\theta$):} Trade flows whose relative contributions are lower than a threshold $\theta$ are ignored, which drastically prunes the search tree.
    \item \textbf{Convergence threshold ($\varepsilon$):} The computation of longer paths stops when their marginal contribution to the accumulated dependency is lower than $\varepsilon$.
    \item \textbf{Maximum path length ($L_{\max}$):} In practice, values between 3 and 5 are used, since longer routes usually have negligible strengths.
    \item \textbf{Use of sparse matrices:} Since most countries only trade with a subset of partners, the use of sparse data structures drastically optimizes memory usage and computation speed.
\end{itemize}

\section{Data and Domestic Supply}\label{sec:data}

\subsection{Trade Flows and the Supply Denominator}
The index is computed from the International Trade and Production Database for Estimation, release 3 (ITPD-E-R03), published by the U.S. International Trade Commission \citep{Borchert2021}. It covers 170 industries (agriculture, mining and energy, manufacturing and services) and 236 economies, and it reports both international and domestic trade, the latter defined as gross production minus exports. For each industry, the trade matrix $\mathbf{X}$ therefore contains domestic trade on its diagonal ($x_{cc}$) and imports off the diagonal, and the supply of country $c$ is
\begin{equation}
S_c = x_{cc} + \sum_{a \neq c} x_{ac}.
\end{equation}
Domestic trade enters the index only through this denominator. Indirect paths exclude the importer and the exporter as intermediaries, so the diagonal never acts as a link in a path, but it determines how much of each country's supply depends on foreign sources, both for the final importer and for every intermediary.

\subsection{Gaps in Domestic Trade}
Domestic trade requires production statistics (UNIDO, FAO, energy and services statistics) that are published with a lag and do not cover all countries and industries. When production is missing, the database omits the domestic row. The share of world imports in country-industry cells with observed domestic trade falls from 54.9\% in 2015 to 45.2\% in 2019, 39.9\% in 2021 and 1.4\% in 2022. In addition, thirteen industries (for instance, live animals, cotton, fishing and some residual services) have no domestic trade for any country or year.

Without correction, the supply denominator of an affected cell collapses to imports. The measure then reflects the concentration of the import basket rather than dependence on total supply: a large producer that imports little would appear fully dependent on its main supplier (for example, crude oil in Saudi Arabia or wine in Italy). The declining coverage would also introduce an artificial upward trend in dependence. The procedure below completes domestic trade while leaving the definition of the index unchanged.

\subsection{Completion Strategy}
The completion follows five principles: (i) observed domestic trade is kept, with the narrow exception of non-credible residuals described in \cref{subsec:credibility}; (ii) the missing values are filled with the same concept, domestic production net of exports, taken from the OECD Inter-Country Input-Output tables \citep{OECDICIO}, which are built from national accounts, complemented by energy statistics for oil producers outside those tables; (iii) the 170 industries and all countries of the ITPD-E are preserved; (iv) values are imputed only where the country imports in the industry, since otherwise the diagonal does not enter any calculation; and (v) every rule was chosen by out-of-sample validation, hiding observed values and measuring the error made when reconstructing them.

\paragraph{Sector correspondence and blocks.} ICIO domestic production of country $c$ in sector $s$ is the sum of the row $(c,s)$ over the columns of country $c$ (intermediate and final uses), which equals output minus exports. The 170 ITPD-E industries are mapped to the 50 ICIO sectors (ISIC Rev.~4), after correcting 22 assignments inherited from the ISIC Rev.~3 classification of the ITPD-E (for example, publishing moves to J58--J60 and rubber and plastics to C22). Industries that share an ICIO sector form a \emph{block} $B$ (40 in total), so that the sum of their domestic trade can be compared with a single ICIO total $Q_{cBt}$.

\paragraph{Calibration.} In country-block-years where all industries are observed, the ratio of aggregated ITPD-E domestic trade to ICIO domestic production has a median of 0.80 in manufacturing (interquartile range 0.57--0.99), 1.03 in mining and energy, and 1.06 in services. Both sources measure the same concept; in manufacturing the ITPD-E is somewhat lower, as industrial statistics do not cover all firms. The median ratio by block, $\kappa_B$, is used in the controls and in the block caps.

\subsection{Imputation Rules}
Let $D_{cit}$ denote domestic trade of country $c$ in industry $i$ (block $B$) and year $t$, $M_{cit}$ imports, $X_{cit}$ exports, and $\rho_{cit} = D_{cit}/Q_{cBt}$. Missing values are filled according to the hierarchy in \cref{tab:diag-rules}.

\begin{table}[!t]
\caption{Imputation rules for missing domestic trade}\label{tab:diag-rules}
\centering
\footnotesize
\begin{tabularx}{\textwidth}{@{}p{1.6cm}p{4.4cm}X@{}}
\toprule
\textbf{Rule} & \textbf{Applies when} & \textbf{Estimate} \\
\midrule
a & The country has credible observed values for the industry within 7 years (before or after) & $\hat D_{cit} = \bar\rho_{ci}\, Q_{cBt}$, where $\bar\rho_{ci}$ averages the ratios of up to three nearest observed years, weighted by the inverse of the distance. Countries outside the ICIO use the ICIO rest-of-world aggregate (a-ROW). \\
b & ICIO country without own history & $\hat D_{cit} = \tilde\rho_{it}\, Q_{cBt}$, where $\tilde\rho_{it}$ is the median ratio of other ICIO countries in $t-1$, $t$ and $t+1$ (at least three distinct countries; otherwise the 2016--2022 median). \\
b-export & Industry without observed domestic trade in any country & $\hat D_{cit} = \kappa_B\, Q_{cBt}\, \bar X_{ci}/\bar X_{cB}$, using the country's 2016--2022 export share of the industry within the block if it is at least 1\%; otherwise, the world export structure of the block. Industries that form a block on their own (e.g.\ fishing) take $\hat D_{cit} = Q_{cBt}$. \\
e & Country outside the ICIO without history, or no reference for rule b & Logit model of the domestic share $\sigma = D/(D+M)$ with industry fixed effects and log trade of the cell and of the country (three-year means), plus the country's own effect when it has at least five observed cells (e+country). $\hat D = \hat\sigma/(1-\hat\sigma)\,\bar M$, with $\hat\sigma$ bounded by the 1st--99th percentiles observed in the industry. \\
en & Crude oil and gas extraction or petroleum refining, country outside the ICIO with physical data & Production from the \emph{Statistical Review of World Energy} \citep{EIStatReview}: refinery throughput, or crude oil plus natural gas output in barrels of oil equivalent, valued at the median unit value of ICIO countries (ICIO output of the sector per barrel, scaled to the ITPD-E), times the country's structural domestic share of production (2016--2022 median of production net of exports). It takes precedence over a-ROW and e. \\
\bottomrule
\end{tabularx}
\end{table}

\paragraph{Energy outside the ICIO.} The rest-of-world aggregate of the ICIO mixes very different economies and is a poor guide for oil producers, whose observed ITPD-E domestic trade is often a low or volatile residual. For crude oil and gas and for refining in countries outside the ICIO, production is therefore taken from sector-specific energy statistics, which cover all major producers over several decades. In a leave-one-country-out test, the rule yields a median error of 4.6 points in the domestic share for refining (5.6 for crude oil and gas). It affects about 0.1\% of world imports but corrects implausible results for large producers (for example, the domestic share of refined products rises from 14\% to 82\% in Iraq and from 7\% to 39\% in Oman in 2022). Observed values below 25\% of the energy-based reference, more than two typical errors below it, are treated as residuals.

\subsection{Credibility of Observed Domestic Trade}\label{subsec:credibility}
Because domestic trade is a residual (production minus exports), it becomes unreliable when declared production barely exceeds exports: the difference is then dominated by measurement error. Examples are pharmaceuticals in Italy, which fall from USD 8.4 billion in 2005 to 0.3 billion in 2015 while exports rise, or aircraft in the United States, with USD 0.6 billion in 2019 between 26.7 billion in 2018 and 52.8 billion in 2020. Used as anchors, such values would propagate near-zero domestic shares to large producers.

Each cell is compared with an export-based reference, the ICIO domestic production of the block allocated according to the industry's weight in the country's exports within the block:
\begin{equation}
R_{cit} = Q_{cBt}\,\frac{X_{cit}}{X_{cBt}}.
\end{equation}
Both export figures come from the ITPD-E, so re-exports cancel out. An observed value with $D_{cit} < \varphi\, R_{cit}$ is treated as a gap and imputed; it is not used as an anchor or in the median ratios. Symmetrically, no imputed value may fall below $\varphi\, R_{ci}$ (computed with the 2016--2022 export share, to avoid breaks across years). We set $\varphi = 0.10$, close to the 10th percentile of $D/R$ among observed cells of large exporters (median 0.70). The criterion affects 4.6\% of observed domestic values in 2009--2022 (1.7--2.5\% of world imports per year) and leaves the out-of-sample error of the rules unchanged (\cref{tab:diag-validation}).

\subsection{Consistency Adjustments}
After imputation, five adjustments ensure coherence across years and with ICIO totals; every cell keeps its previous value and rule for traceability:
\begin{enumerate}[leftmargin=*,itemsep=2pt]
    \item \textbf{Safeguard for rule a.} If rule a moves the domestic share more than 30 percentage points away from the nearest observed value, or yields a value more than five times larger or smaller, the whole country-industry series uses the level of the nearest observation.
    \item \textbf{Splicing between rules.} When a series alternates years with and without own history, the latter are rescaled to link with the nearest year based on own history, preserving their own dynamics.
    \item \textbf{Magnitude cap.} No value from rule b exceeds 1.5 times the largest observed value of the industry in any country (when at least ten countries are observed), except for the three largest world exporters.
    \item \textbf{Block caps.} No imputed industry exceeds $1.5\,\kappa_B\, Q_{cBt}$; if observed plus imputed values of the block exceed $2\,\kappa_B\, Q_{cBt}$, the imputed part without own history is scaled down by a constant factor per country-block (minimum 25\%).
    \item \textbf{Export-coherence floor.} No imputed value falls below $0.10\,R_{ci}$.
\end{enumerate}

\subsection{Validation}
\Cref{tab:diag-validation} reports the out-of-sample error of the selected rules. Observed domestic values for 2016--2021 are hidden and reconstructed, and the error is measured on the domestic share of supply, $\sigma = D/(D+M)$, weighted by imports, which is the component of the index affected by the diagonal. Rules based on cross-country medians exclude the country being predicted, and the logit model is validated by country groups not used in the estimation.

\begin{table}[!t]
\caption{Out-of-sample validation of the imputation rules (2016--2021)}\label{tab:diag-validation}
\centering
\footnotesize
\begin{tabularx}{\textwidth}{@{}Xrr@{}}
\toprule
\textbf{Rule} & \textbf{Cells} & \textbf{Error in $\sigma$ (pp)} \\
\midrule
a: own history, three anchors, inverse distance & 23,081 & 2.39 \\
\quad with credible anchors and export-coherence floor & 22,303 & 2.42 \\
\quad carrying forward the last level (benchmark) & 23,081 & 3.18 \\
a-ROW: own history, rest-of-world dynamics & 10,316 & 4.31 \\
b: median ratio of other countries & 23,642 & 9.65 \\
\quad agriculture / manufacturing / mining and energy / services & & 16.1 / 12.4 / 11.2 / 2.3 \\
b-export: country export structure (industries without any observation) & 23,371 & 12.37 \\
\quad world export structure (discarded) & 20,166 & 22.09 \\
e+country: logit model with country effect (outside the ICIO) & 9,583 & 14.07 \\
Median domestic share of the industry (discarded) & 9,702 & 29.61 \\
\bottomrule
\end{tabularx}
\par\smallskip\parbox{\linewidth}{\footnotesize Error: mean absolute error in the domestic share of supply, weighted by imports, in percentage points.}
\end{table}

Own history remains preferable to the cross-country rule even when the nearest observation is six or seven years old (error of 5.0 and 8.7 points, against 10.3 and 16.1 for rule b), which supports the seven-year window.

\subsection{Quality Controls}
Each imputed cell passes seven controls with thresholds fixed in advance: validity; consistency of the block sum with $\kappa_B\,Q_{cBt}$; ratio to the nearest own observation; year-on-year continuity; change in the domestic share relative to the nearest observation; plausibility of the domestic share against the 1st--99th percentiles observed in the industry; and magnitude relative to the largest observed value. A deviation counts only if the alternative reference value would move the domestic share by at least five percentage points. Between 0.4\% and 0.9\% of imputed imports fall in the red category in each year; all such cells carry a documented decision (for example, the largest world producer of an industry, or a genuine change in ICIO production) and none remains pending.

\subsection{Integration in the Index}
Domestic trade is incorporated at a single point, right after reading the ITPD-E and before building the matrices; in the EU version, before aggregating the 27 member states, so that EU domestic supply includes each member's domestic trade and intra-EU trade. Diagonal rows are added only for countries and industries already present in the data, so the set of countries and industries does not change; non-credible residuals are replaced in place. The relevance threshold $\theta$ of the computation (0.5\% of the column total) is kept unchanged; with complete domestic supply, that total includes domestic production, as it already did for observed cells.

\subsection{Updates and Revisions}\label{subsec:revisions}
The domestic value of a year draws on neighbouring years: rule a anchors on the closest observations before or after it, the rule b ratio uses a three-year window, and a few structural parameters (export weights, the energy domestic share and $\kappa_B$) are estimated over a calibration window, 2016--2022. To keep the published series stable, the calibration window is part of the methodology version and is not re-estimated when a year is added. When data for year $t$ become available, the index is computed for $t$ and recomputed for $t-1$, which stops being the edge of the series; earlier years keep their published values. As with GDP and other economic statistics, figures for the two latest years are therefore provisional. A simulation that adds 2022 to a series ending in 2021 supports this rule: with fixed parameters, about 1\% of world imports change in 2016--2020, with a negligible effect on the import-weighted domestic share (at most 0.02 points), whereas 4\% change in 2021. The whole series is recomputed and released as a new version only when the methodology changes or when a new edition of the ITPD-E or the ICIO revises past years. A year enters the index only once the ICIO covers it.

\subsection{Limitations}
Rule b, the least precise, is applied to roughly a third of world imports, with errors of 16 points in agriculture and 12 in manufacturing but only 2 in services. Industries without any observed domestic trade rely on structural assumptions (error of 12 points, and up to 27 points for marginal exporters). Countries that assemble exports with imported inputs from the same industry (processing trade) show high dependence on their suppliers even when they are large exporters; this reflects industry-level trade data rather than the imputation. Finally, the latest ICIO years include OECD estimates, whose effect is limited by the safeguard of rule a.


\section{Centrality Analysis and Global Hubs (Global Hub Score)}

The identification of critical intermediaries is one of the major contributions of the EIED. Countries that act as "connectors" in trade networks can amplify crises.

\subsection{Intermediation Metrics}
For each country $k$, we define two metrics:
\begin{itemize}
    \item \textbf{Intermediation Frequency ($\phi_k$):} The number of significant trade routes in which country $k$ participates as an intermediary.
    \item \textbf{Intermediation Strength ($\psi_k$):} Weighted sum of the strength of those routes, assigning greater importance to intermediaries closer to the final importer.
\end{itemize}

\subsection{Composite Centrality Index (GHS)}
The \textit{Global Hub Score} ($C_k$) integrates both dimensions:
\begin{equation}
C_k = \alpha \cdot \frac{\phi_k}{\max \phi} + (1-\alpha) \cdot \frac{\psi_k}{\max \psi}
\end{equation}
Empirically, $\alpha = 0.4$ offers a robust balance, allowing countries to be classified according to their structural power as global "hubs" sustaining the world trade architecture.

\section{Disruption Simulation and Resilience Analysis}

To holistically assess vulnerability, we implement counterfactual scenario simulations. This involves the hypothetical removal of key nodes (such as the supply of a critical sector in a dominant country) to quantify the structural impact on the rest of the network. We measure:
\begin{itemize}
    \item \textbf{Connectivity impact:} Degree to which other countries become isolated or suffer critical supply shortages.
    \item \textbf{Centrality reconfiguration:} How structural power shifts toward other nodes after a disruption.
    \item \textbf{Hidden vulnerabilities:} Identification of dependencies that only manifest after the fall of a previous intermediary.
\end{itemize}

\section{Conclusions and Implications for Economic Policy}

The main contribution of the Elcano's Index on Economic Dependence (EIED) is the exposure of "hidden" vulnerabilities. By decomposing dependency into its direct and indirect components, the design of balanced policy responses is facilitated.

Direct dependencies can be mitigated through traditional diversification of immediate suppliers. However, indirect dependencies require an "open strategic autonomy" strategy that includes fostering domestic capacities in critical supply links and forming strategic alliances with partners presenting complementary dependency profiles.

This indicator provides a solid quantitative basis for "structural surveillance", allowing systemic risks to be anticipated before they turn into open trade crises, and guiding the design of more resilient supply chains in an uncertain geopolitical environment.

\bibliographystyle{apalike}
\bibliography{references.bib}

\end{document}
